\newpage
\subsection{Sicherheit der Lieferkette}
\textbf{Anforderungen:} NIS2 fordert die Sicherheit der Lieferkette einschließlich sicherheitsbezogener Aspekte der Beziehungen zu unmittelbaren externen Dienstleistern. Ziel ist es, Verfügbarkeitsrisiken und Supply-Chain-Angriffe zu verhindern.

Die Sicherheit der Lieferkette umfasst zwei Aspekte:
\begin{itemize}
    \item \textbf{Verfügbarkeitsrisiken:} Der Ausfall eines Dienstleisters gefährdet die eigene Geschäftstätigkeit (z.\,B. Cloud-Anbieter ist nicht erreichbar)
    \item \textbf{Sicherheitsrisiken:} Ein kompromittierter Dienstleister wird für Angriffe auf die eigene Organisation genutzt (Supply-Chain-Attacken)
\end{itemize}

Voraussetzung für die Bewertung ist eine vollständige Übersicht über alle kritischen Dienstleister und ihre Leistungen; ohne diese Transparenz lassen sich weder Verfügbarkeits- noch Sicherheitsrisiken bewerten. Die systematische Erfassung und Dokumentation aller NIS2-relevanten Dienstleister ist daher unverzichtbar. Empfehlenswert ist, die Lieferkette nach einem Business Continuity Management bzw. einer Business Impact Analyse anzugehen, da diese die Abhängigkeiten kritischer Geschäftsprozesse von externen Dienstleistern aufzeigt.

\textbf{Kategorisierung der Dienstleister:} Bewerten Sie Ihre Dienstleister nach beiden Risikoaspekten:
\begin{itemize}
    \item \textbf{Verfügbarkeitskritische Dienstleister:} Deren Ausfall Ihre Produktion oder Verfügbarkeit direkt beeinträchtigt (Cloud-Anbieter, Rechenzentren, kritische SaaS-Anwendungen)
    \item \textbf{Sicherheitskritische Dienstleister:} Die Zugriff auf Ihre IT-Infrastruktur haben oder sensible Daten verarbeiten (Managed Service Provider, IT-Support, Softwareentwickler)
    \item \textbf{Beide Kategorien:} Viele Dienstleister fallen in beide Kategorien und erfordern besondere Aufmerksamkeit
\end{itemize}

NIS2 fordert explizit, Supply-Chain-Attacken zu verhindern, bei denen Angreifer zunächst einen vertrauenswürdigen Drittanbieter kompromittieren, um über diesen „Umweg" an das eigentliche Ziel zu gelangen. Potenzielle Angriffe und ihre Auswirkungen müssen bewertet werden.
\newpage
\textbf{Herausforderung der unklaren Anforderungen:} Welche Sicherheitsmaßnahmen von Dienstleistern zu fordern sind, lässt das NIS2-Umsetzungsgesetz bewusst offen. Das BSI verlangt lediglich „möglichst konkrete Anforderungen an die Cybersicherheit" und überlässt die Definition den Unternehmen. Explizit genannt werden nur:
\begin{itemize}
    \item Risikomanagementmaßnahmen
    \item Maßnahmen zur Bewältigung von Cybersicherheitsvorfällen
    \item Patchmanagement
    \item Prinzipien wie Security by Design oder Security by Default
\end{itemize}

Diese Unsicherheit erfordert einen pragmatischen Ansatz. Zertifikate wie ISO 27001, BSI IT-Grundschutz, SOC 2, TISAX etc. bieten eine gute Grundlage und sind als Nachweis gut dokumentierbar. Die individuelle Bewertung eines Dienstleisters kann aber auch ergeben, dass ein Zertifikat nicht ausreicht.
\newpage
Nach Auslegung des BSI sollten Dienstleister NIS2-betroffener Unternehmen selbst mindestens die Anforderungen aus Artikel 21 der NIS2-Richtlinie umsetzen. Als Orientierung empfiehlt das BSI das 30-seitige Dokument „Best-Practice-Empfehlungen für Anforderungen an Lieferanten zur Gewährleistung der Informationssicherheit in Kritischen Infrastrukturen" der UP KRITIS mit folgenden Themen:
\begin{itemize}
    \item Verantwortung des Auftraggebers und Auftragnehmers
    \item Vulnerability-Management
    \item Patch-Management
    \item Systemhärtung (Minimale Prinzipien, Sicherheitsstandards, Passwörter etc.)
    \item Fernzugang für Drittanbieter
    \item Anforderungen an Softwareentwicklungsprozesse
    \item Einsatz von Kryptografie
    \item Sicherheitsanforderungen an den IT-Betrieb
    \begin{itemize}
        \item Informationssicherheitsprozesse/ISMS
        \item Zugriffschutz und Rechtevergabe
        \item Asset-Management
        \item Personalsicherheit
        \item Physische Sicherheit
        \item Netzwerksicherheit und operationelle Sicherheit
        \item Security-Incident-Management
    \end{itemize}
    \item Dokumentation
    \item Informationspflicht über sicherheitsrelevante Vorfälle
    \item Nicht-technische Sicherheit
    \begin{itemize}
        \item Organisation der Informationssicherheit
        \item Asset-Management
        \item Human-Resources-Security
        \item Audits
    \end{itemize}
\end{itemize}

Die Themen decken sich größtenteils mit den NIS2-Vorgaben. Lesen Sie das Dokument, berücksichtigen Sie es bei der Bewertung von Dienstleistern und dokumentieren Sie, dass Sie sich daran orientiert haben.
\newpage
\textbf{Praktisches Vorgehen bei Dienstleister-Bewertungen:}

\textbf{1. Erfassung:} Definieren Sie einen strukturierten Prozess zur Erfassung aller NIS2-relevanten Dienstleister. Nicht alle Dienstleister müssen erfasst werden, aber klare Bewertungskriterien (z.\,B. Zugriff auf IT-Systeme, Verarbeitung kritischer Daten, Einfluss auf Geschäftsprozesse) sichern Kontinuität und verhindern, dass kritische Dienstleister übersehen werden.

Implementieren Sie folgende Teilprozesse:
\begin{itemize}
    \item \textbf{Regelmäßige Überprüfung:} Die Dienstleisterliste sollte mindestens jährlich auf Korrektheit, Vollständigkeit und Aktualität überprüft werden
    \item \textbf{Onboarding neuer Dienstleister:} Neue Dienstleister werden vor Vertragsabschluss systematisch nach den definierten Kriterien bewertet und bei Relevanz erfasst
    \item \textbf{Offboarding:} Nach Beendigung der Geschäftsbeziehung werden Dienstleister zeitnah aus dem System entfernt und zugehörige Zugriffsrechte widerrufen
\end{itemize}

Dokumentieren Sie Bewertungskriterien und Verantwortlichkeiten für jeden Teilprozess.

\textbf{2. Priorisierung:} Konzentrieren Sie sich zunächst auf Dienstleister, die Teil Ihrer IT-Infrastruktur sind oder direkten Einfluss darauf haben, etwa Cloud-Anbieter oder Rechenzentren.

\textbf{3. Gestaffelte Anforderungen:} Die Anforderungstiefe richtet sich nach der Kritikalität des Dienstleisters:
\begin{itemize}
    \item \textbf{Basis-Level:} Nachweis gültiger Sicherheitszertifikate (ISO 27001, SOC 2 etc.)
    \item \textbf{Erweitert:} Zusätzliche Dokumentation zu spezifischen Sicherheitsmaßnahmen
    \item \textbf{Detailliert:} Umfassende Sicherheitsbewertung für höchstkritische Dienstleister
\end{itemize}

\textbf{4. Kommunikation mit Dienstleistern:} Weisen Sie bei Anfragen auf die NIS2-Richtlinie hin. Unternehmen außerhalb der EU kennen NIS2 möglicherweise nicht, obwohl sie durch ihre Dienste in der EU betroffen sein könnten. Fragen Sie gezielt nach:
\begin{itemize}
    \item Gültigen Sicherheitszertifikaten und deren Erneuerungszyklen
    \item IT-Risikomanagement- und Patchmanagement-Konzepten
    \item Spezifischen Sicherheitsmaßnahmen für Ihre Verbindung (VPN-Verschlüsselung, SSH-Schlüssel-Management etc.)
    \item NIS2-Betroffenheit des Dienstleisters selbst
\end{itemize}
\newpage
\textbf{5. Dokumentation und Risikobewertung:} Alle Antworten sind zu dokumentieren. Antwortet ein Dienstleister nicht oder sind seine Maßnahmen unzureichend, bedeutet das nicht automatisch eine Vertragsbeendigung, ist aber in der Risikobewertung zu berücksichtigen:
\begin{itemize}
    \item Bei zu hohem Risiko: Dienstleister-Wechsel erwägen
    \item Alternative: Dem Dienstleister eine Frist zur Nachbesserung setzen
    \item Für künftige Verträge: Cybersicherheitsanforderungen vertraglich festhalten
\end{itemize}

Dokumentieren Sie Bewertungslogik, Prozesse und Entscheidungskriterien sorgfältig. Das hilft bei künftigen Audits und belegt die Nachvollziehbarkeit Ihres Vorgehens trotz der unklaren regulatorischen Anforderungen.
